\chapter*{Реферат}
\addcontentsline{toc}{chapter}{Реферат}
Пояснительная записка содержит \pageref{end_of_document}~стр., 2~рис., 7~табл. Количество использованных источников 50. Количество приложений 0.\par
Суть исследования состоит в разработке и экспериментальной проверке прототипа модуля семантического сопоставления вакансий и резюме для системы автоматизированного подбора персонала. Модуль использует модель раздельного кодирования, обученную на исторических парах текст вакансии и текст резюме с метками кадрового решения, и возвращает числовую оценку соответствия кандидата требованиям вакансии.\par
\textbf{Ключевые слова:} система автоматизированного подбора персонала (Applicant Tracking System, ATS), подбор персонала, резюме, описание вакансии, семантическое сопоставление вакансии и резюме, ранжирование кандидатов, модель раздельного кодирования вакансии и резюме (bi-encoder), модель векторных представлений предложений (Sentence-BERT), косинусная мера сходства, фреймворк языка Python под названием FastAPI, тестовый контур системы автоматизированного подбора персонала под названием Reqcore.\par
В первом разделе анализируются предметная область подбора персонала, системы автоматизированного подбора персонала, методы обработки естественного языка и подходы к ранжированию кандидатов. По итогам раздела формулируется постановка задачи разработки модуля семантического сопоставления вакансий и резюме.\par
Во втором разделе рассматриваются исходные данные вакансий и резюме, требования к их очистке, разметке и преобразованию в обучающие пары. Также описываются программный интерфейс сервиса ранжирования и его взаимодействие с тестовым контуром системы автоматизированного подбора персонала Reqcore.\par
В третьем разделе описываются реализация конвейера машинного обучения, дообучение модели RuBERT Tiny 2, сохранение артефактов модели и устройство класса ранжирования ATSRanker. Отдельно фиксируется сервис на базе FastAPI, который предоставляет маршруты индексации, ранжирования и оценки одной пары текстов.\par
В четвертом разделе проводится экспериментальная оценка качества, проверка интеграции с Reqcore и расчет временного и стоимостного эффекта внедрения. Сравниваются три стратегии ранжирования на отложенных данных: BM25, перекрестная модель совместного кодирования и модель раздельного кодирования.\par
В приложениях дополнительные материалы не приводятся.\par
